\documentclass[a4paper,10pt]{article}
\usepackage[utf8]{inputenc}
\usepackage[T1]{fontenc}
\usepackage[french]{babel}
\usepackage{enumitem}
\usepackage[left=1.4cm,right=1.4cm,top=1cm,bottom=1cm]{geometry}
\usepackage{hyperref}
\usepackage{xcolor}
\usepackage[sfdefault,light]{roboto}
\renewcommand{\familydefault}{\sfdefault}

\pagestyle{empty}
\setlength{\parindent}{0pt}
\definecolor{linkcolor}{HTML}{1A73E8}
\hypersetup{colorlinks=true, urlcolor=linkcolor}

% Separator used in the contact line and in titles
\newcommand{\sep}{\textbar\ }
% Section heading: bold caps with a rule under it
\newcommand{\cvsection}[1]{\vspace{0.55em}\noindent{\textbf{\MakeUppercase{#1}}}\par\vspace{-0.4em}\noindent\rule{\linewidth}{0.4pt}\par\vspace{0.25em}}
% Left text, right text on one line
\newcommand{\cvrow}[2]{\noindent #1\hfill #2\par}
% Bulleted list
\newenvironment{cvlist}{\begin{itemize}[leftmargin=1.1em,itemsep=0pt,topsep=1pt,parsep=0pt,partopsep=0pt,label=\textbullet]}{\end{itemize}\vspace{0.3em}}

\begin{document}
\fontsize{10}{12}\selectfont

%=========== EN-TÊTE ===========
\begin{center}
  {\fontsize{22}{24}\selectfont\textbf{Marsel Bakashov}}\\[0.3em]
  +1 (514) 664-9320 \sep \href{mailto:m.bakashov11@gmail.com}{m.bakashov11@gmail.com} \sep \href{https://www.linkedin.com/in/marsel-bakashov-813643213/}{LinkedIn} \sep \href{https://github.com/Mbaka11}{github.com/Mbaka11} \sep \href{https://marselbak.netlify.app/}{Portfolio}
\end{center}
\vspace{-0.6em}

%=========== RÉSUMÉ ===========
\cvsection{Résumé}
Étudiant à la maîtrise en \textbf{apprentissage automatique} (Mila), avec des stages en \textbf{systèmes multi-agents}, \textbf{apprentissage par renforcement} et \textbf{recherche}. J'aime transformer l'apprentissage automatique en outils réellement utilisés, surtout là où les idées doivent résister aux vraies données, contraintes et utilisateurs. J'apprends en construisant, et j'aime enseigner et organiser. Je cherche un stage où prendre un problème en charge de bout en bout, avec de fortes équipes.

%=========== ÉDUCATION ===========
\cvsection{Éducation}
\cvrow{\textbf{Université de Montréal / McGill -- Mila}}{Montréal, Canada}
\cvrow{\textit{Maîtrise pro en apprentissage automatique, GPA 4.0 / 4.0}}{\textit{Jan. 2026 -- Jan. 2028}}
\vspace{0.2em}
\cvrow{\textbf{Polytechnique Montréal}}{Montréal, Canada}
\cvrow{\textit{Baccalauréat en génie logiciel, GPA 3.7 / 4.0}}{\textit{Sept. 2021 -- Déc. 2025}}
\vspace{0.2em}
\cvrow{\textbf{Korea University College of Engineering}}{Séoul, Corée du Sud}
\cvrow{\textit{Baccalauréat en génie informatique (échange académique)}}{\textit{Fév. 2025 -- Juin 2025}}

%=========== EXPÉRIENCE ===========
\cvsection{Expérience}
\cvrow{\textbf{Stagiaire en science des données et ML} \sep \textit{McKesson}}{Mai 2026 -- Déc. 2026}
\begin{cvlist}
  \item \textbf{Développé} des agents IA utilisés par \textbf{plus de 100 personnes} pour analyser les données d'affaires et générer des synthèses, économisant environ \textbf{10 h par semaine par utilisateur} (estimé).
  \item \textbf{Conçu} à grande échelle des workflows multi-agents orchestrés : des dizaines d'agents dotés d'outils et de compétences dédiés traitent les données à la chaîne en se passant leurs résultats.
  \item \textbf{Traité} des données de santé sensibles (PHI) via des API, en priorisant confidentialité et sécurité.
\end{cvlist}

\cvrow{\textbf{Stagiaire en science des données et RL} \sep \textit{Intact Assurance}}{Sept. 2025 -- Avr. 2026}
\begin{cvlist}
  \item \textbf{Développé} des politiques de tarification par apprentissage par renforcement qui améliorent la rentabilité à long terme en simulation, sous \textbf{plus de 10 contraintes} réglementaires et d'affaires.
  \item \textbf{Validé} les politiques en simulation sur des \textbf{millions d'enregistrements} (rentabilité, volatilité, CVaR).
\end{cvlist}

\cvrow{\textbf{Stagiaire en science des données} \sep \textit{Radio-Canada / CBC}}{Mai 2024 -- Sept. 2024}
\begin{cvlist}
  \item \textbf{Conçu} un système d'indexation basé sur l'IA unifiant les sources de contenu internes pour la recherche sémantique et lexicale sur \textbf{plus de 15 M d'éléments}, utilisé par \textbf{plus de 100 employés}.
\end{cvlist}

\cvrow{\textbf{Développeur full-stack stagiaire} \sep \textit{Stingray}}{Sept. 2023 -- Déc. 2023}
\begin{cvlist}
  \item \textbf{Développé} des fonctionnalités full-stack en équipe agile : analytique, tests, CI et revues de code.
\end{cvlist}

\cvrow{\textbf{Assistant d'enseignement, mathématiques discrètes} \sep \textit{Polytechnique Montréal}}{Jan. 2024 -- Août 2024}
\begin{cvlist}
  \item \textbf{Animé} des révisions, corrigé des devoirs et soutenu \textbf{80+ étudiants} en logique et graphes.
\end{cvlist}

\cvrow{\textbf{Directeur des commandites} \sep \textit{PolyAI Hackathon}}{Jan. 2024 -- Mai 2024}
\begin{cvlist}
  \item \textbf{Obtenu} \textbf{22\,000\,\$} en financement et dirigé une équipe de commandites avec réunions.
\end{cvlist}

%=========== PROJETS ===========
\cvsection{Projets}
\cvrow{\textbf{ApplyPilot} \sep \textit{JavaScript --- Node.js --- Extension Chrome --- LLM}}{\href{https://github.com/Mbaka11/ApplyPilot}{GitHub}}
\begin{cvlist}
  \item \textbf{Conçu} une extension Chrome et un serveur local qui analysent l'offre, choisissent le CV, rédigent la lettre avec un LLM et remplissent les formulaires, avec \textbf{plus de 380 tests} ; l'utilisateur envoie.
\end{cvlist}

\cvrow{\textbf{GravityGuyML} \sep \textit{Python --- RL profond --- Gymnasium --- TensorFlow --- Pygame}}{\href{https://github.com/Mbaka11/Gravity-Guy-ML}{GitHub}}
\begin{cvlist}
  \item \textbf{Conçu} un environnement Pygame compatible avec Gymnasium, intégrant une génération procédurale et des agents d'apprentissage profond (DQN, PPO, SAC, A2C).
\end{cvlist}

\cvrow{\textbf{Prévision d'épidémies à partir de données climatiques} \sep \textit{Python --- scikit-learn --- XGBoost}}{\href{https://github.com/Mbaka11/20251R0136COSE47101}{GitHub}}
\begin{cvlist}
  \item \textbf{Prédit} 11 diagnostics à partir de symptômes et de données météo (environ 5\,000 cas) avec KNN, SVM et XGBoost, pour un \textbf{F1 macro de 0,98} et un \textbf{AUC-ROC de 0,999}.
\end{cvlist}

%=========== COMPÉTENCES ===========
\cvsection{Compétences}
\textbf{Langages}~: Python, SQL, C++, Java, Bash, R\\
\textbf{ML / Données}~: PyTorch, Scikit-learn, TensorFlow, Hugging Face, MLflow, apprentissage par renforcement, agents IA\\
\textbf{Outils}~: Azure, Databricks, Snowflake, Power BI, Airflow, Docker, Git\\
\textbf{Langues}~: anglais, français, russe, espagnol, coréen

\end{document}
